\section{双目标前沿、组件消融与75片排程}\label{sec:pareto-evidence}

\subsection{两目标的冲突与可认证前沿}

两片算例S1含同一Pjob的两片$2\times2$产品、四片辅助片，清洁阈值为2。采用该工况的配对与双腔时序子模型，在82.03~s内重新生成图~\ref{fig:pareto-s1}的七份排程。逐点重放并与独立数值证书比对，在$2\times10^{-5}$容差内命中七个全局非支配点，含三个非支持点。最快回片为331~s，对应$B=424/3$。最低$B=42$对应519~s。当上限为$B\le100$时，两片分别与辅助片配对，最短工期为514~s；只从线性加权和的支持点中选解则需519~s。该例的$R=0$，两目标的冲突来自交付与路线等待。专用重算的HV为0.378339，$\mathrm{IGD}^{+}=2.45\times10^{-10}$。通用领域搜索在40~s内每个种子命中一个认证点，种子1延长至300~s后命中两个点。
\begin{figure}[!htbp]
\centering\includegraphics[width=.92\linewidth]{pareto_s1}
\caption{两片算例S1在两种配对方式下的完整前沿。两图分别标注原始坐标，三枚菱形为线性加权和无法取得的非支持点。}\label{fig:pareto-s1}
\end{figure}

四片标称混合工况M4中，发现了413~s交付的排程。已证最低$B$为54，达到该值至少需要476~s，已验证排程将对应的最短工期限定在$[476,497]$~s。因此追求最低$B$至少要比413~s的排程多等63~s。在已访问的35种整数结构中，开放连续时刻后有一段位于这些结构合并边界上的有效线段，其工期区间为$(420.5,423)$~s。该区间的结论限于这35种结构。

\subsection{等预算对照与扰动}

五个工况各比较六种方法，每个方法使用三个种子。S1、M4、H4、H8和H16的产品组成（$4\times1$片数，$2\times2$片数）依次为$(0,2)$、$(2,2)$、$(3,1)$、$(4,4)$和$(8,8)$。同工况、同种子的各臂共享可行起点。NSGA-II采用非支配排序与拥挤距离选择\citep{deb2002}，MOEA/D采用切比雪夫分解与邻近权重子问题替换\citep{zhangli2007}。两种基线以Python实现并共用设备解码器。$\varepsilon$-SCIP直接在原模型中扫描目标上限，原两阶段法保持工期优先。表~\ref{tab:pareto-small}给出各方法返回集合的超体积（HV）中位数。对归一化的双目标计算该集合支配至参考点$(1.1,1.1)$的面积，数值越大表示覆盖越广。S1按认证前沿极值定标，其他工况按共同起点$(A_0,B_0)$取$[0,1.5A_0]\times[0,2B_0]$。HV在同一工况内比较。
\begin{table}[!htbp]\centering\small
\caption{五个工况的三种子HV中位数与单次搜索时限}\label{tab:pareto-small}
\begin{tabular}{lrrrrrrr}\toprule
工况 & 时限/s & 领域搜索 & 加选割 & $\varepsilon$-SCIP & NSGA-II & MOEA/D & 两阶段 \\\midrule
@@SMALL_TABLE@@
\bottomrule\end{tabular}
\end{table}

领域搜索在两种四片混合工况的HV中位数均最高。八片、十六片工况由MOEA/D取得最高中位数。用同一随机种子配对比较，M4开启配对下界、在已有下界时采用领域邻域、允许改变对内取放顺序，对应的HV差值中位数分别为$+0.00372$、$+0.01710$和$+0.06650$，三项改动在三个种子中均提高HV。自适应选择上限的差值为$-0.02102$，接入冻结选割策略的差值为$-0.00330$。在12对固定上限、共享起点的4~s原模型子问题中，下界对其中6对排除了至少一种配对数量，12对的条件工期均相同。

M4另将邻域释放比例从0.30分别改为0.15和0.50，两个配对种子的HV差值中位数分别为$-0.03107$和$-0.02767$。把修复预算比例从0.22改为0.12和0.32，对应变化为$-0.00734$和$+0.00105$。在八片节拍实验R8中，151份归档排程有46份产生非零节拍偏差。对同批候选在相同工期上限下分别按0、12.5、25和50重算节拍权重，五个有可选排程的种子与工期上限组合均选择了相同的$R$值。这些数据区分了节拍项实际激活与改变权重后决策发生变化两种情况。

\subsection{75片混合批次的可行前沿}

采用默认配置D，37片执行$4\times1$，38片执行$2\times2$。LP1至LP4的两类产品数量分别为$(10,9)$、$(9,10)$、$(9,10)$和$(9,9)$。两个Pjob共用CH2和CH3，真空辅助库存为8片，清洁阈值为3，时域为10000~s。历史构造得到的可行起点为$(A_0,B_0)=(8630,3932)$。四种方法由同一起点出发，各运行三个种子，每次限时90~s。共享起点的外部准备与完整原模型的事后复核单独计时，搜索时限包含结构调整、重定时和设备事件复验。固定结构对照仅改变动作时刻及$B$上限，在预设网格的扫描区间耗尽后提前停止。

每次运行按预设规则提交至多七点，先保留两个端点，再依次选择与已保留集合的归一化距离最大的点。表~\ref{tab:pareto75}的HV仅使用当次提交点。最后两列分别对各次提交集合的最小$A$与最小$B$取中位数，两列可以对应不同的排程。
\begin{table}[!htbp]\centering\small
\caption{75片混合批次的三种子结果。HV范围为三次运行的最小值至最大值。}\label{tab:pareto75}
\begin{tabular}{lrrrr}\toprule
方法 & HV中位数 & HV范围 & 最小$A$/s & 最小$B$/加权秒 \\\midrule
@@LARGE_TABLE@@
\bottomrule\end{tabular}
\end{table}

领域搜索的HV中位数为@@DHV@@，MOEA/D为@@MHV@@。第二个种子中MOEA/D取得@@BESTSEED@@，高于同种子的领域搜索。固定生产块和资源次序后，单靠重定时仍能在$(8630,3932)$与$(8852,3722)$之间找到取舍。开放结构后的三个搜索方法均改善了固定结构法的中位HV。75片领域搜索在本组所用上限上未触发配对下界筛除，表中差异不能归因于该下界或学习选割。

图~\ref{fig:pareto75}汇集全部12次运行在限时内发现的非支配解，含单次提交时被七点上限省略的内部点。按两目标再次去支配后共得@@NFRONT@@个实际坐标。在最短工期端，$(A,B)=(@@AMIN@@,@@BATAMIN@@)$；在最低等待与节拍指标端，$(A,B)=(@@AATBMIN@@,@@BMIN@@)$。从前者换到后者，批次晚交付@@DELTA_A@@~s，$B$减少@@DELTA_B@@加权秒。若只在本次候选中选择，要求$B\le2000$时最早可在@@A2000@@~s交付。该值可与两端点一起用于选择交期。
\begin{table}[!htbp]\centering\small
\caption{75片共同起点与按交付要求选取的三份排程。$A$、$L_{\rm sum}$、$L_{\max}$和$R$的单位为秒，$B$为加权秒。}\label{tab:pareto75-policies}
\begin{tabular}{lrrrrr}\toprule
排程选择 & $A$ & $B$ & $L_{\rm sum}$ & $L_{\max}$ & $R$ \\\midrule
共同起点 & 8630 & 3932 & 3312 & 62 & 0 \\
发现的最短工期 & 7567 & 2599 & 1884 & 49 & 9 \\
$B\le2000$时最早交付 & 7639 & 1951 & 1776 & 17.5 & 0 \\
发现的最低$B$ & 8045 & 1782 & 1612 & 17 & 0 \\\bottomrule
\end{tabular}
\end{table}
\FloatBarrier
\begin{figure}[!htbp]
\centering\includegraphics[width=.96\linewidth]{pareto_75}
\caption{75片混合批次的经验前沿。左图叠加各方法提交点、共同起点与全部发现记录的合并非支配点，右图放大最短工期端。每个非支配点均有独立核验的完整可行排程。}\label{fig:pareto75}
\end{figure}

两端点的排程均含九个可计量的节拍位置。沿最短工期端到最低$B$端，总等待由1884降至1612~s，最大单项等待由49降至17~s，节拍偏差由9降至0~s。三项分别解释$B$下降量中的272、320和225加权秒。

75片的12次运行共提交@@NSUB@@份排程。对提交集合及补全后的合并前沿去重，共有@@NCERT@@份排程逐片重放，并向原模型的6\,764\,011个变量赋值、逐一检查24\,271\,470条约束。按事件复算的$A$、$B$与原模型赋值相符。此前小规模主比较连同消融、参数扰动和长预算实验共164组，699份最终排程均已核对。小规模的完整性证书仅适用于S1。75片的网格、有限结构邻域与搜索预算限定了本图的覆盖范围。对完整可行域尚无排除其他非支配点的证书。
